Governments in centralized administrative systems have to rely on delegation and incentive contracts to implement national goals \citep{williamson1985economic}. However, asymmetric information between upper-level and local governments, combined with the fact that local officials face promotion incentives tied to performance across multiple tasks such as promoting economic growth while meeting environmental protection requirements, can lead to deviations from policy objectives \citep{holmstrom1991multitask, chen2018career}. A common response from the center is to strengthen monitoring and performance-based assessments in an effort to narrow information gaps and realign local incentives embedded in these contracts. Yet when monitoring remains incomplete and is observable on some dimensions but not others, tighter supervision raises the possibility that distortions may be shifted rather than eliminated.

When the center tightens supervision over dimensions it can readily observe, does this induce distortions along dimensions that remain imperfectly monitored? This paper addresses this question by studying China’s tightening of cropland protection policy and its implications for local officials’ land allocation choices. Specifically, I examine whether stricter enforcement of cropland quantity targets reshapes the trade-off local governments face between allocating land to urban expansion, which supports economic growth, and allocating land to cropland to satisfy binding protection requirements.

Local governments in China have long been evaluated primarily on economic performance, particularly economic growth \citep{fan2025land}. To support industrialization and urban expansion, local officials have strong incentives to convert land for construction use, making land allocation a central instrument of local development policy. At the same time, concerns over food security prompted the central government to introduce the “red line” cropland protection policy in the early 2000s, which mandates a minimum national cropland area.\footnote{The red line policy stipulates that total arable land should remain no less than 120 million hectares, reflecting national concerns over grain self-sufficiency.} Local governments thus operate in a multitask environment that requires balancing economic development with cropland protection.

Prior to 2017, however, cropland protection was weakly enforced despite the formal existence of quantitative targets. This changed with the 2017 reform of the cropland protection responsibility assessment system, under which compliance with cropland protection targets became a binding component of official evaluations and promotion criteria. While the reform aimed to strengthen the joint management of cropland quantity and productive capacity, in practice enforcement focused almost exclusively on observable cropland quantity. Monitoring of cropland quality and subsequent utilization remained limited. This asymmetry raises a key concern: when quantitative targets become binding under incomplete monitoring, local governments may comply by expanding cropland area through the reclamation of marginal or low-quality land, generating distortions in land allocation.

In this paper, I use satellite-based raster data to identify cropland and cultivation intensity. The unit of observation is a $500 \times 500$ meter grid cell. Aggregating raster measures to the grid level allows me to track changes in land allocation and land use at a spatial scale relevant for local policy implementation, while avoiding reliance on administrative statistics that may be subject to reporting biases.

I leverage cross-county exogenous variation in cropland protection pressure to identify the effects of tighter cropland protection on land allocation and cultivation outcomes. Cropland protection pressure captures the extent to which local governments are constrained in meeting cropland protection requirements given ongoing land conversion demands. It arises from institutional features of China’s cropland occupation--replenishment system, which allow replenishment obligations to be met outside the converting county. In cases where a county or district lacks sufficient land, replenishment is, in principle, arranged through cross-county adjustment within the same prefecture, with available land across the prefecture serving as a shared reservoir.

As a result, counties differ in their exposure to binding cropland protection constraints depending on their relative endowment of available land within the prefecture. Prefecture-level cropland obligations act as a common policy shock, while pre-existing county-level land endowments determine differential exposure. This institutional structure generates plausibly exogenous variation in cropland protection pressure, which the analysis exploits to study how stricter enforcement reshapes land-use decisions and cultivation outcomes.

This study finds that tighter enforcement of cropland protection affects land-use decisions along both the extensive and intensive margins of land use. On the extensive margin, counties facing greater cropland protection pressure respond more strongly to the post-2017 assessments, with grid cells becoming significantly more likely to be converted into cropland. On the intensive margin, however, cropland reclaimed after the policy in high-pressure counties tends to be of lower soil suitability and is less likely to be actively cultivated, exhibiting higher fallow rates in subsequent years. Further analysis points to multitask incentives as a key mechanism: when promotion incentives are strong, and officials face heightened pressure to balance economic growth with cropland protection, compliance with binding quantity targets is achieved more through the reclamation of low-quality land. Although the policy tightening was intended to realign local implementation with central objectives, the evidence suggests that land misallocation and inefficiency can arise in practice when monitoring of land quality and utilization remains incomplete.

This study contributes to several strands of the literature. First, it is closely related to the multitasking principal–agent framework \citep{holmstrom1991multitask, baker1992incentive}. In contrast to a large empirical literature that examines how governments trade off opposing tasks in the Chinese context \citep{hong2025order, fan2025land, chen2018career, he2020watering}, this paper shifts attention to within-task distortions under imperfect monitoring. It provides the first empirical evidence that stricter performance measures may backfire when oversight is incomplete. Moreover, the existing multitasking literature has largely focused on the balance between environmental and economic objectives. To the best of my knowledge, this paper is the first to examine multitasking incentives in land-use regulation, thereby contributing to the broader literature on land regulation \citep{cai2017build, hilber2016impact, besley2000land}.

This study also contributes to the literature on misallocation and agricultural productivity \citep{adamopoulos2022misallocation, restuccia2017land, lagakos2013selection}. While \cite{adamopoulos2022misallocation} links misallocation in China to distortionary policies, their analysis focuses on productivity losses arising from misallocation across agricultural production factors. By contrast, this paper shows how imperfect performance evaluation systems can generate land misallocation across land-use types, highlighting a distinct channel through which policy-induced distortions affect agricultural efficiency.

Finally, this paper contributes to the literature on the political economy of the Chinese economy. A large body of work emphasizes that the central government frequently relies on top-down planning and performance-based mandates to implement nationwide policies, yet implementation outcomes often deviate from central objectives due to policy design constraints, information asymmetries, or strategic behavior by local governments \citep{wang2025policy, FismanWangYe2025_StrategicCaptureOfMonitors, axbard2024informed, henderson2022political}. Existing studies primarily focus on the logic of policy design, the role of information transmission, or the emergence of collusion. This paper adds to this strand of literature by showing how implementation distortions can arise even in the absence of collusion or explicit manipulation. When local governments face binding multitask incentives under imperfect monitoring, policy compliance along observable dimensions can naturally generate resource misallocation and productivity losses along less observable ones.

The remainder of the paper is organized as follows. Section~\ref{sec:background} describes the institutional background in China, with a focus on cropland protection policy. Section~\ref{sec:data} presents the data and explains the construction of the cropland protection pressure measure. Section~\ref{sec:fact} documents a set of motivating facts. Section~\ref{sec:concept} introduces a conceptual framework that guides the empirical analysis. Section~\ref{sec:empirical} presents the empirical strategy and main results. Section~\ref{sec:conclude} concludes.
